
ABSTRACT 
Recently, there has been a persistent rise in the use of deep learning models to identify and classify plant leaf diseases. Nevertheless, gap exists in the particular methods that can enhance the performance of these models. This project aims to develop and train MobileNetV2 model plant leaf images to detect several disease of a particular plant and classify them with reliable and efficient techniques to increase the reliability of the MobileNetV2 model. The methodology leverages on an experiment to evaluate different approaches for improving the performance of the MobileNetV2 model to detect disease and classify them in plant leaf images. The findings showed the efficiency of various implementation techniques in enhancing the effectiveness of the plant leaf disease detection and classification model.

CHAPTER ONE
INTRODUCTION


CHAPTER TWO
LITERATURE REVIEW




CHAPTER THREE
MATERIALS AND METHODS

	References
	
	Abdu, A. M., Mokji, M. M., \& Sheikh, U. U. (2020). Machine learning for plant disease  
	detection: An investigative comparison between support vector machine and deep learning. IAES International Journal of Artificial Intelligence, 9(4), 670–683. https://doi.org/10.11591/ijai.v9.i4.pp670-683
	
	Ahila Priyadharshini, R., Arivazhagan, S., Arun, M., \& Mirnalini, A. (2019). Maize leaf 
	disease classification using deep convolutional neural networks. Neural Computing and Applications, 31, 8887–8895.
	
	Albattah, W., Nawaz, M., Javed, A., et al. (2022). A novel deep learning method for detection
	and classification of plant diseases. Complex \& Intelligent Systems, 8, 507–524.
	Ang, K. L. M., \& Seng, J. K. P. (2021). Big data and machine learning with hyperspectral information in agriculture. IEEE Access, 9, 36699–36718. https://doi.org/10.1109/ACCESS.2021.3051196
	Bansal, P., Kumar, R., \& Kumar, S. (2021). Disease detection in apple leaves using deep convolutional neural network. Agriculture, 11(7), 617.
	Barure, V., et al. (2020). In revolutionizing crop disease detection with computational deep learning: A comprehensive review.
	Bukumira, M., Antonijevic, M., Jovanovic, D., Zivkovic, M., Mladenovic, D., \& Kunjadic, G. (2022). Carrot grading system using computer vision feature parameters and a cascaded graph convolutional neural network. Journal of Electronic Imaging, 31(6), 061815.
	Chen, C. J., Huang, Y. Y., Li, Y. S., Chen, Y. C., Chang, C. Y., \& Huang, Y. M. (2021). Identification of fruit tree pests with deep learning on embedded drone to achieve accurate pesticide spraying. IEEE Access, 9, 21986–21997.
	Chen, L., Lin, S., Lu, X., Cao, D., Wu, H., Guo, C., ... \& Wang, F. Y. (2021). Deep neural network-based vehicle and pedestrian detection for autonomous driving: A survey. IEEE Transactions on Intelligent Transportation Systems, 22(6), 3234–3246. https://doi.org/10.1109/TITS.2020.2993926
	Ekhuemelo, C., \& Otor, E. O. (2021). Assessment of fungi associated with the leaf spot disease of fluted pumpkin (Telfairia occidentalis Hook F.) and their management using botanicals. Nigerian Journal of Mycology, 13, 31–43.
	Fu, L., Li, S., Sun, Y., Mu, Y., Hu, T., \& Gong, H. (2022). Lightweight-convolutional neural network for apple leaf disease identification. Frontiers in Plant Science, 13, 831219.
	Ghosh, A., \& Roy, P. (2022). An automated model for leaf image-based plant recognition: An optimal feature-based machine learning approach. Innovations in Systems and Software Engineering, 20(4), 705–718. https://doi.org/10.1007/s11334-022-00440-y
	Hosny, K. M., El-Hady, W. M., Samy, F. M., Vrochidou, E., \& Papakostas, G. A. (2023). Multi-class classification of plant leaf diseases using feature fusion of deep convolutional neural network and local binary pattern. IEEE Access, 11, 62307–62317.
	Jiang, P., Chen, Y., Liu, B., He, D., \& Liang, C. (2019). Real-time detection of apple leaf diseases using deep learning approach based on improved convolutional neural networks. IEEE Access, 7, 59069–59080.
	Khalid, M., Sarfraz, M. S., Iqbal, U., Aftab, M. U., Niedbała, G., \& Rauf, H. T. (2023). Real-time plant health detection using deep convolutional neural networks. Agriculture, 13(2), 510. https://doi.org/10.3390/agriculture13020510
	Khakimov, A., Salakhutdinov, I., Omolikov, A., \& Utaganov, S. (2022). Traditional and current-prospective methods of agricultural plant diseases detection: A review. IOP Conference Series: Earth and Environmental Science, 951(1), 012002. https://doi.org/10.1088/1755-1315/951/1/012002
	Khan, A. I., Quadri, S. M. K., Banday, S., \& Shah, J. L. (2022). Deep diagnosis: A real-time apple leaf disease detection system based on deep learning. Computers and Electronics in Agriculture, 198, 107093.
	Liu, J., \& Wang, X. (2021). Plant diseases and pests detection based on deep learning: A review. Plant Methods, 17, 22. https://doi.org/10.1186/s13007-021-00722-9
	Massi, I. E., Mostafa, Y. E.-S., Yassa, E., \& Mammass, D. (2020). Combination of multiple classifiers for automatic recognition of diseases and damages on plant leaves. Signal, Image and Video Processing. https://doi.org/10.1007/s11760-020-01797-y
	Mgbemena, I. (2024). Interesting facts you should know about fluted pumpkin (Ugu). AgroNigeria, January.
	Nwaneto, C. B., \& Yinka-Banjo, C. (2024). Harnessing deep learning algorithms for early plant disease detection: A comparative study and evaluation between SSD (MobileNet_v2 and MobileNet_v3) and CNN model. Ife Journal of Science, 26(3).
	Sadhu, S., \& Kaligotla, V. V. S. K. (2023). Improving the accuracy of plant leaf disease detection and classification in images of plant leaves: By exploring various techniques with the MobileNetV2 model (Bachelor’s thesis, Blekinge Institute of Technology).
	Sun, J., Yang, Y., He, X., \& Wang, X. (2020). Northern maize leaf blight detection under complex field environment based on deep learning. IEEE Access, 8, 33679–33688. https://doi.org/10.1109/ACCESS.2020.2973658
	Waheed, A., Goyal, M., Gupta, D., Khanna, A., Hassanien, A. E., \& Pandey, H. M. (2020). An optimized dense convolutional neural network model for disease recognition and classification in corn leaf. Computers and Electronics in Agriculture, 175, 105456.
	Wei, W. U., Yang, T. L., Rui, L. I., Chen, C. H. E. N., Tao, L. I. U., Zhou, K., Sun, C. M., Li, C. Y., Zhu, X. K., \& Guo, W. S. (2020). Detection and enumeration of wheat grains based on a deep learning method under various scenarios and scales. Journal of Integrative Agriculture, 19(8), 1998–2008. https://doi.org/10.1016/S2095-3119(19)62803-0
	Wu, H., Wiesner-Hanks, T., Stewart, E. L., DeChant, C., Kaczmar, N., Gore, M. A., \& Lipson, H. (2019). Autonomous detection of plant disease symptoms directly from aerial imagery. Plant Phenome Journal, 2(1), 1–9.
	Xie, X., Ma, Y., Liu, B., He, J., Li, S., \& Wang, H. (2020). A deep-learning-based real-time detector for grape leaf diseases using improved convolutional neural networks. Frontiers in Plant Science, 11, 751. https://doi.org/10.3389/fpls.2020.00751
	Zamani, A. S., Anand, L., Rane, K. P., Prabhu, P., Buttar, A. M., \& Dugbakie, B. N. (2022). Performance of machine learning and image processing in plant leaf disease detection. Journal of Food Quality, 2022, Article ID 1598796.
	Zhang, C., Wang, J., Yan, T., et al. (2023). An instance-based deep transfer learning method for quality identification of Longjing tea from multiple geographical origins. Complex \& Intelligent Systems, 9, 3409–3428.
	Zhou, C., Zhou, S., Xing, J., \& Song, J. (2021). Tomato leaf disease identification by restructured deep residual dense network. IEEE Access, 9, 28822–28832. https://doi.org/10.1109/ACCESS.2021.3058947
